\documentclass[11pt]{article}
\usepackage[T1]{fontenc}
\usepackage[utf8]{inputenc}
\usepackage{lmodern}
\usepackage[margin=1.05in]{geometry}
\usepackage{amsmath,amssymb,amsthm,mathtools}
\usepackage{booktabs,array}
\usepackage{xcolor}
\usepackage[colorlinks=true,linkcolor=blue!45!black,citecolor=blue!45!black,urlcolor=blue!45!black]{hyperref}
\numberwithin{equation}{section}
\newtheorem{theorem}{Theorem}[section]
\newtheorem{proposition}[theorem]{Proposition}
\newtheorem{lemma}[theorem]{Lemma}
\theoremstyle{definition}
\newtheorem{definition}[theorem]{Definition}
\theoremstyle{remark}
\newtheorem{remark}[theorem]{Remark}
\newcommand{\F}{\mathbb F}
\newcommand{\R}{\mathbb R}
\newcommand{\Z}{\mathbb Z}
\title{A kissing configuration of 12\,277 points in dimension 19}
\author{Claude (AI model by Anthropic)\\ \small run and submitted by Aleksandr Bersenev}
\date{October 14, 2026\\[3pt]\footnotesize Produced autonomously by Claude (Opus 5.5) agents in Claude Code; run and submitted by Aleksandr Bersenev.\\ \footnotesize Not an official Anthropic publication. Not peer reviewed.}
\hypersetup{pdftitle={A kissing configuration of 12277 points in dimension 19},pdfauthor={Claude (AI model by Anthropic); submitted by Aleksandr Bersenev}}
\begin{document}
\maketitle
\begin{abstract}
We exhibit $12\,277$ points on the unit sphere $S^{18}\subset\R^{19}$ with pairwise angular
distances at least $60^\circ$. Hence the kissing number in dimension $19$ satisfies
$\tau_{19}\ge 12\,277$, one more than the bound $12\,276$ of K.~Li (October 2026). The
configuration keeps Li's $10\,484$-point base, which comes from the odd-sign configuration of Cohn and
Li. It restores one more of the $184$ base points that Li deletes, and it replaces Li's $1792$ added
points by $1792$ re-optimised points. These form $112$ orbits under the order-$16$ group of sign
changes that fixes the restored point. All coordinates are integers. Every pair is checked in exact
arithmetic by two independent programs. The statement is also proved in Lean~4 by kernel evaluation
without \texttt{native\_decide}: the inner products among base points are bounded by a formalised
structural lemma, and all other inner products are reduced by the symmetry group to finitely many
integer inequalities. We also report numerical evidence that restoring a whole further orbit of
eight deleted points, as asked by Li, is impossible in this family: the best minimax cosine we
reach is $\tfrac12+2.0\cdot10^{-5}$.
\end{abstract}

\section{Introduction}
The kissing number $\tau_n$ is the largest number of non-overlapping unit balls in $\R^n$ that touch
a fixed unit ball. Equivalently, it is the largest size of a set of unit vectors with pairwise inner
products at most $1/2$. It is known only for $n\in\{1,2,3,4,8,24\}$; see Cohn's table~\cite{CohnTable}
for current bounds and references. For $17\le n\le 23$ the classical lower bounds come from the
Leech lattice cross sections $\Lambda_n$~\cite{LS71}, which give $\tau_{19}\ge 10\,668$. Cohn and
Li~\cite{CL24} changed the parity of the sign patterns on the octads and added sign vectors of a
code; this gave $\tau_{19}\ge 11\,692$. Ho~\cite{Ho26} found a larger code and reached $11\,948$.
In October 2026 K.~Li~\cite{Li26,Li26b} obtained $\tau_{19}\ge 12\,276$. That construction cuts the
$20$-dimensional Cohn--Li configuration by a hyperplane, deletes $184$ base points and adds
$1792$ reweighted sign vectors with tails. Li lists as open whether ``a second orbit of eight
deleted points [can] be kept, under $\Gamma$ or under a smaller group with a larger
table''~\cite[\S6]{Li26b}. On October 15, 2026, Cohn's table~\cite{CohnTable} lists $12\,276$ as
the best lower bound and $24\,417$ as the best upper bound~\cite{LdL24}.

\begin{theorem}\label{thm:main}
There are $12\,277$ unit vectors in $\R^{19}$ with pairwise inner products at most $1/2$. Hence
$\tau_{19}\ge 12\,277$.
\end{theorem}

The configuration of Theorem~\ref{thm:main} answers a one-point version of Li's question positively:
one more deleted point can be kept if the symmetry group is shrunk to its stabiliser and the added
points are re-optimised. For a whole extra orbit of eight points we found only configurations with
largest cosine $\tfrac12+2.0\cdot10^{-5}$ (Section~\ref{sec:limits}). To our knowledge, as of
October 15, 2026, Theorem~\ref{thm:main} is the best lower bound for $\tau_{19}$. We checked
Cohn's table, the four public repositories that worked on dimensions $17$--$31$ in September and
October 2026~\cite{Li26,Li26b,Kra26,Long26}, the preprint~\cite{YZW26} and a web search.

\section{Notation and the base}\label{sec:base}
We follow~\cite{Li26b}. Coordinates of $\R^{20}$ are numbered $1,\dots,20$. Let $E\subseteq\F_2^{20}$
be the $[20,12,4]$ code spanned by the twelve words of~\cite[\S3.1]{Li26b}, which is the Golay code
punctured at four coordinates. Its dual contains $130$ \emph{octads} (words of weight $8$), and two
octads meet in at most $4$ coordinates. The \emph{odd configuration} $B_{20}$ of Cohn and
Li~\cite{CL24} consists of the $760$ roots $\pm2e_i\pm2e_j$ together with, for every octad $o$, the
$128$ vectors with entries $\pm1$ on $o$, zeros elsewhere and an odd number of minus signs. All
these vectors have squared norm $8$.

\begin{lemma}[Cohn--Li]\label{lem:CL}
Distinct $b,b'\in B_{20}$ satisfy $\langle b,b'\rangle\le4$.
\end{lemma}
\begin{proof}
The two cases that need an argument are the following. Two vectors on the same octad have odd sign
patterns, so they differ in an even and positive number of places. Two vectors on different octads
overlap in at most four coordinates. The remaining cases are immediate.
\end{proof}

Let $T=\{1,4,7,9\}$, and let $H=\{1,\dots,20\}\setminus T$ (the \emph{head}). The points of $B_{20}$
orthogonal to $\mathbf 1_T$ number $10\,668$. Among them are $192$ points on the four octads
containing $T$ whose sign pattern on $T$ has two plus and two minus signs. Li deletes $184$ of these
and keeps the eight points $P$ on the octad $T\cup\{3,6,8,19\}$ with pattern $(+,+,-,-)$ on $(1,4,7,9)$.
The map $v\mapsto(v_H, Rv_T/2)$, with
\[R=\begin{pmatrix}1&1&-1&-1\\1&-1&1&-1\\1&-1&-1&1\end{pmatrix},\]
is an isometry from $\mathbf 1_T^\perp$ onto $\R^{19}$~\cite[(2)]{Li26b}. We work in
$\mathbf 1_T^\perp\subset\R^{20}$ throughout, so ``a kissing configuration in $\mathbf 1_T^\perp$''
means one in $\R^{19}$.

Li's sign group $\Gamma$ consists of the $128$ sign changes $S_g$ on the coordinates of words
$g$ of a $7$-dimensional subcode of $E$ missing $T$~\cite[\S3.2]{Li26b}. Every $S_g$ maps the base
onto itself. Under $\Gamma$ the $184$ deleted points form $23$ orbits of size $8$. In Li's
configuration each such orbit conflicts with exactly one orbit of added points, at cosines between
$0.5286$ and $0.5525$.

\section{The configuration}\label{sec:config}
Let
\[ q_0 = (-1,-1,+1,+1 \text{ on } 1,4,7,9;\ -1,-1,-1,+1 \text{ on } 10,15,17,18;\ 0 \text{ elsewhere}). \]
This is one of Li's deleted points: it lies on the octad $T\cup\{10,15,17,18\}$, has five minus signs
and is orthogonal to $\mathbf 1_T$. Let $\mathcal H\le\Gamma$ be its stabiliser. It has order $16$
and consists of the $S_g$ with $g\cap\{10,15,17,18\}=\emptyset$. The configuration $\mathcal C$ has
three parts.
\begin{itemize}
\item \textbf{Base:} the $10\,668$ points of $B_{20}\cap\mathbf 1_T^\perp$ minus the $192$ points
described above, plus $P$ and $q_0$. This gives $10\,485=492+9\,993$ points ($492$ roots and $9\,993$ octad
vectors). It is $\mathcal H$-invariant.
\item \textbf{New points:} the $\mathcal H$-orbits $\{S_g r:\ S_g\in\mathcal H\}$ of $112$ integer vectors
$r_1,\dots,r_{112}\in\mathbf 1_T^\perp\cap\Z^{20}$. Their entries have absolute value at most
$1.16\cdot10^{8}$, and they are listed in the data file \texttt{reps112\_R20.txt}. Every $r_i$ has all
$16$ head coordinates nonzero, so each orbit has exactly $16$ points, and the $1792$ new points have
$1792$ distinct head sign patterns. These patterns are the head patterns of $1792$ words $c\in E$
with $|c\cap T|\in\{1,2,3\}$, chosen as one bipartition class in each of the $112$ cubes of such
words. (The cubes are the cosets of the span of the five tetrads.)
\end{itemize}
Thus $|\mathcal C| = 10\,485+1\,792=12\,277$.

\begin{proposition}\label{prop:check}
All $12\,277$ vectors of $\mathcal C$ are distinct, nonzero and orthogonal to $\mathbf 1_T$. Every
pair $x\ne y$ satisfies $\langle x,y\rangle\le0$ or $4\langle x,y\rangle^2\le|x|^2|y|^2$. Moreover,
for every pair involving a new point, $\tfrac14-\cos^2\angle(x,y)\ge 4.37\cdot10^{-6}$ whenever the
inner product is positive.
\end{proposition}
Theorem~\ref{thm:main} follows by normalising and applying the isometry of Section~\ref{sec:base}.
For the proof of Proposition~\ref{prop:check} we use the symmetry reduction of~\cite[\S3.3]{Li26b}.
Base pairs are covered by Lemma~\ref{lem:CL}. A pair $(b, S_g r_i)$ is carried by the isometry $S_g$
to $(S_gb, r_i)$, and $S_gb$ lies in the base. A pair $(S_gr_i,S_hr_j)$ is carried to
$(r_i,S_{g+h}r_j)$. It therefore suffices to check two lists:
\begin{itemize}
\item the $112\times10\,485$ products $\langle r_i,b\rangle$ with $b$ in the base;
\item the $112\cdot112\cdot16$ products $\langle r_i,S_gr_j\rangle$ with $(g,j)\neq(0,i)$.
\end{itemize}
We verified the whole statement in three independent ways (Section~\ref{sec:verif}).

\paragraph{How the configuration was found.} We parametrised Li's family by representatives modulo
a sign group and minimised the largest cosine by sequential linear programming on the
symmetry-reduced list of pairs. Keeping the base, $P$ and the full $\Gamma$-orbit of $q_0$, and
optimising over all $2^7$ $\Gamma$-invariant choices of kept words, the best minimax value is
$\tfrac12+2.0\cdot10^{-5}$. From that configuration we dropped seven of the eight points of the
orbit. We then ran a gradient method (Adam on a squared-hinge penalty) on all $1792$ added points
without symmetry, which reached $\tfrac12+1.2\cdot10^{-5}$. We averaged the result over
$\mathcal H$ and polished it by $\mathcal H$-symmetric sequential linear programming, reaching
$\tfrac12-4.37\cdot10^{-6}$. Finally we rounded at scale $10^8$. The search played no part in the
proof.

\section{Verification}\label{sec:verif}
\paragraph{Two independent exact programs.} The first program computes all $75\,356\,226$ cosines
in floating point and re-checks the $2\,390\,370$ pairs with cosine above $\tfrac12-10^{-6}$ in exact
integer arithmetic. The second computes exact $64$-bit Gram blocks, which is safe because
$|x_k|<2^{28}$, and checks $4\langle x,y\rangle^2\le|x|^2|y|^2$ with unbounded integers for all
$26\,085\,106$ pairs with positive inner product. Both pass. Exactly $2\,390\,370$ pairs touch, all
of them base--base. Both programs also pass on the $19$-coordinate version of the data.

\paragraph{Lean.} The certificate uses core Lean~4.35, with no Mathlib needed. Kernel caches accumulate
across declarations, so to bound memory the certificate is split into $33$ modules, compiled with
\texttt{lean -o} in dependency order (\texttt{lean/mods2/}). They are: generic lemmas; the data;
$14$ modules of cheap checks; and $16$ modules with $14$ kernel checks each, taking $6$--$10$ minutes
each. The last module contains the main theorem. The same content concatenated into one file is
\texttt{lean/Main.lean}. The main theorem is
\begin{verbatim}
theorem K19.kissing19 :
  exists L : List (List Int), L.length = 12277 /\ L.Nodup /\
    (forall v in L, v.length = 20 /\ tsumT v = 0 /\ 0 < dot v v) /\
    L.Pairwise Compat
\end{verbatim}
(written with ASCII words for $\exists$, $\wedge$, $\forall$, $\in$).
Here \texttt{dot} is the integer inner product of lists and \texttt{tsumT v} is $v_0+v_3+v_6+v_8$
(the $0$-based indices of $T$). \texttt{Compat x y} is
$\langle x,y\rangle\le0\vee 4\langle x,y\rangle^2\le\langle x,x\rangle\langle y,y\rangle$. The list
$L$ is computed inside Lean from the data: $128$ octad records (octad mask and list of sign masks),
$492$ roots, $16$ sign masks and $112$ representatives. The proof uses no \texttt{sorry},
\texttt{axiom}, \texttt{admit} or \texttt{native\_decide}. It has four parts.
\begin{enumerate}
\item \emph{Base pairs: a formalised structural proof.} The bit-level identity
  $|o\cap(m\oplus m')|+2|o\cap m\cap m'|=|o\cap m|+|o\cap m'|$ shows that two odd octad vectors on one
  octad differ in an even, positive number of coordinates, so their inner product is at most $4$.
  For different octads the inner product is at most $|o\cap o'|$, and the $128^2$ octad
  intersections are checked by \texttt{decide}. Roots are handled by a two-term formula.
\item \emph{Group closure.} $\mathcal H$ is closed under composition, and every octad block and the
  root list are closed under the sign changes. Both are checked by \texttt{decide}.
\item \emph{Reduced inequalities.} $224$ theorems, two per representative, are checked by
  \texttt{decide +kernel}. They cover the $10\,485$ base products, evaluated through a proved
  $8$-term formula for octad vectors, and the $1\,792$ products with sign-changed representatives.
\item \emph{Assembly.} The isometry lemmas
  $\langle S_gx,S_gy\rangle=\langle x,y\rangle$ and $S_gS_h=S_{g\oplus h}$ give
  \texttt{L.Pairwise Compat}. Distinctness follows because no nonzero vector is compatible with
  itself.
\end{enumerate}
The command \verb|#print axioms K19.kissing19| reports \texttt{[propext, Classical.choice,
Quot.sound]}. Resource figures are in \texttt{lean/lean\_output.txt}.

\section{Limits and open questions}\label{sec:limits}
\begin{itemize}
\item \emph{A whole extra orbit.} We ran the minimax optimisation for every combination of three
  choices: one of the $2^7$ $\Gamma$-invariant choices of kept words; one of seven candidate deleted
  orbits; and Li's group $\Gamma$ or one of the other $15$ admissible index-$2$ subgroups of the
  sign group of $E$ fixing $T$ (sampled). The best value is $\tfrac12+2.0\cdot10^{-5}$. The minima
  take a handful of discrete values. Passing to subgroups of index $2$, $4$ or $8$ at the optimum
  does not lower it, and neither does letting $P$, the extra orbit or base orbits move. This is
  numerical evidence, not a proof, that Li's question has a negative answer for full orbits.
\item \emph{Two extra points.} Two points of the same orbit (same stabiliser $\mathcal H$) reach
  $\tfrac12+1.75\cdot10^{-5}$ by $\mathcal H$-symmetric optimisation. Whether $\tau_{19}\ge12\,278$
  can be reached this way is open.
\item \emph{Dimension 21.} Li's line-weighted points satisfy a rule (his Lemma~4.3) implying that the
  excluded words in each coset of the line code form a dominating set of the Cayley graph of
  $\F_2^{10}$ with respect to the $21$ lines of $\mathrm{PG}(2,4)$. Hence this family gives at most
  $4(1024-\gamma)$ points, where $\gamma$ is the domination number. Tabu search finds dominating sets
  of size $64$ immediately and none of size $63$. This suggests $\gamma=64$, in which case Li's
  $3840$ points would be optimal for line weights. We have no proof.
\end{itemize}

\paragraph{Use of AI.} This work, including the search, the verification code, the Lean
formalisation and this text, was carried out by Claude (Anthropic).

\begin{thebibliography}{99}
\bibitem{CohnTable} H.~Cohn, \emph{Table of kissing number bounds}, \url{https://cohn.mit.edu/kissing-numbers}, accessed October 15, 2026; data set \url{https://hdl.handle.net/1721.1/153312}.
\bibitem{CL24} H.~Cohn and A.~Li, \emph{Improved kissing numbers in seventeen through twenty-one dimensions}, preprint, 2024, arXiv:2411.04916.
\bibitem{Ho26} B.~S.~Ho, \emph{A new lower bound for the kissing number in 19 dimensions}, preprint, 2026, arXiv:2603.10425.
\bibitem{Li26} K.~Li, \emph{D19--D21 kissing configurations}, \url{https://github.com/Felpix-Studios/D19-D21-Kissing}, October 2026.
\bibitem{Li26b} K.~Li, \emph{Improved kissing numbers in nineteen and twenty-one dimensions}, preprint in \url{https://github.com/Felpix-Studios/kissing-number}, October 2026.
\bibitem{Kra26} A.~Kravatskiy, \emph{Kissing numbers: new lower bounds in dimensions 18 through 96}, \url{https://github.com/alexlegeartis/KissingNumbers}, 2026.
\bibitem{Long26} G.~Long, \emph{Kissing configurations in $\R^{21}$}, \url{https://github.com/griffinwork40/D21-Kissing}, October 2026.
\bibitem{YZW26} S.~Yang et al., \emph{Large-scale autonomous discovery of kissing number constructions}, preprint, 2026, arXiv:2609.35051.
\bibitem{LS71} J.~Leech and N.~J.~A.~Sloane, \emph{Sphere packings and error-correcting codes}, Canad.\ J.\ Math.\ 23 (1971), 718--745.
\bibitem{LdL24} N.~Leijenhorst and D.~de Laat, \emph{Solving clustered low-rank semidefinite programs arising from polynomial optimization}, Math.\ Program.\ Comput.\ 16 (2024), 503--534.
\end{thebibliography}
\end{document}
